\exnumber{04}
\extitle{UX V - Ex 5 - Wireframes}
\exfiles{ex04/Figma file, named Name\_FirstName\_DeliverableName\_Date\_VersionNumber}
\exauthorize{Those named in the exercise}
\makeheaderfiles

\textbf{Instructions:}

Pick a product (a record, a song or an album) from your music label website.

Based on the new taxonomy (content categorization), show the different paths a user has to take to reach the product with a series of wireframes.

\textbf{Tips:}

\begin{itemize}
\item Drawing a user journey first might be helpful before rushing into wireframes
\item Drawing paper prototypes before going digital might save you some time ;)
\item Several paths are encouraged
\end{itemize}

Vocabulary

\begin{itemize}
\item Low-Fidelity: Paper Prototypes - Example here \url{https://miro.medium.com/max/1400/1*5nUfqqA2gjdAYHagjbHA5w.jpeg}
\item Mid-Fidelity: Wireframes - black and white, only the main elements, on digital prototype tools - Example here \url{https://miro.medium.com/max/1400/1*eGSfaxwYbxUFomYbyO6GWw.png}
\item High-Fidelity: Wireframes with color, styles, graphical details, and micro-interactions - Example here \url{https://miro.medium.com/max/1400/1*Xn0HSKAvhr4TZzC9lN5udw.gif}
\end{itemize}

\textbf{List of websites:}

\begin{itemize}
\item Universal Music Latin America \url{https://www.universalmusica.com/}
\item Delta Records \url{https://deltarecords.net/}
\item Zoho Music \url{https://www.zohomusic.com/}
\item Luaka Bop \url{https://www.luakabop.com/}
\item Rimas Music \url{https://rimasmusic.com/}
\item More... as long as it is a South American music label ;-)
\end{itemize}

\textbf{Deliverables:}

\begin{itemize}
\item Wireframes (6 to 12 screens) from the home page to the product page, on Figma
\item Upload the Figma file on Github, labeled as follows: "Name\_FirstName\_DeliverableName\_Date\_VersionNumber". For example, it can be labeled as "Doe\_John\_MidFiPrototype\_05242024\_V1".
\item Add a title within the file.
\item Pay attention to the global aspect of the file. It must be clear, simple and easy to read. You can get inspiration on canva.com but don't overload your design with too much details!
\end{itemize}

\textbf{Resources:}

\begin{itemize}
\item How to start sketching and wireframing \url{https://uxdesign.cc/how-to-start-sketching-and-wireframing-84a821f092e2}
\item Why should product teams use wireframes more often? \url{https://uxdesign.cc/why-should-product-teams-use-wireframes-more-often-60e34a2bc55}
\item The 5-pass reduction wireframing: A minimalist UX technique \url{https://uxdesign.cc/the-5-pass-reduction-wireframing-a-minimalist-ux-technique-aa415aad9ce2?sk=a218ac4db9d1049249bd381fbc5edeee}
\item The art of designing good wireframes \url{https://uxdesign.cc/the-art-of-designing-good-wireframes-8a2a9c10a95b}
\end{itemize}

\textbf{Recommended tools:}

\begin{itemize}
\item Figma
\end{itemize}

